\hypertarget{appendix-m-the-algorithm-compendium-the-masters-toolkit}{%
\section{Appendix M: THE ALGORITHM COMPENDIUM (The Master's
Toolkit)}\label{appendix-m-the-algorithm-compendium-the-masters-toolkit}}

Welcome to the Master's Toolkit. Most C++ developers write
\passthrough{\lstinline!for!} loops. Gods use
\passthrough{\lstinline!<algorithm>!}. Why? Because the algorithms in
the STL are already optimized, exception-safe, and carry semantic
meaning. When you see \passthrough{\lstinline!std::partition!}, you
immediately know what the code is doing. When you see a 20-line
\passthrough{\lstinline!for!} loop, you have to play computer in your
head to figure it out.

The following is a comprehensive, ``Godhood-level'' breakdown of the
110+ functions available in \passthrough{\lstinline!<algorithm>!},
\passthrough{\lstinline!<numeric>!}, and
\passthrough{\lstinline!<memory>!}. This is not just a list; it is a
tactical guide to hardware-aware, expressive, and high-performance C++
programming.

\begin{center}\rule{0.5\linewidth}{0.5pt}\end{center}

\hypertarget{stdall_of}{%
\subsubsection{\texorpdfstring{1.
\texttt{std::all\_of}}{1. std::all\_of}}\label{stdall_of}}

\begin{itemize}
\tightlist
\item
  \textbf{Analogy}: The ``Strict Bouncer''. If even one person in the
  line doesn't have an ID, nobody gets in.
\item
  \textbf{When to use it}: When you need to verify that a property holds
  for an entire collection (e.g., ``Are all these packets valid?'').
\item
  \textbf{Complexity}: \(O(n)\).
\item
  \textbf{Common Pitfall}: Forgetting that an empty range returns
  \passthrough{\lstinline!true!} (Vacuous Truth).
\item
  \textbf{Hardware Sympathy}: Short-circuits immediately. If the first
  element fails, the CPU doesn't even fetch the rest of the array into
  the cache.
\item
  \textbf{Example}:
  \passthrough{\lstinline!cpp     std::vector<int> v = \{2, 4, 6, 8\};     bool all\_even = std::all\_of(v.begin(), v.end(), [](int i)\{ return i \% 2 == 0; \});!}
\end{itemize}

\hypertarget{stdany_of}{%
\subsubsection{\texorpdfstring{2.
\texttt{std::any\_of}}{2. std::any\_of}}\label{stdany_of}}

\begin{itemize}
\tightlist
\item
  \textbf{Analogy}: The ``Optimist''. As long as one person has a
  ticket, the party is a success.
\item
  \textbf{When to use it}: To check if at least one element satisfies a
  condition (e.g., ``Is there any corrupted data in this block?'').
\item
  \textbf{Complexity}: \(O(n)\).
\item
  \textbf{Common Pitfall}: Returns \passthrough{\lstinline!false!} for
  empty ranges.
\item
  \textbf{Hardware Sympathy}: High branch misprediction potential if the
  matching element is near the middle of a large range.
\item
  \textbf{Example}:
  \passthrough{\lstinline!cpp     bool has\_negative = std::any\_of(v.begin(), v.end(), [](int i)\{ return i < 0; \});!}
\end{itemize}

\hypertarget{stdnone_of}{%
\subsubsection{\texorpdfstring{3.
\texttt{std::none\_of}}{3. std::none\_of}}\label{stdnone_of}}

\begin{itemize}
\tightlist
\item
  \textbf{Analogy}: The ``Clean Slate''. Ensuring there are no spiders
  in the room.
\item
  \textbf{When to use it}: To verify that no elements satisfy a negative
  condition.
\item
  \textbf{Complexity}: \(O(n)\).
\item
  \textbf{Hardware Sympathy}: Effectively
  \passthrough{\lstinline"!any\_of"}. Compilers often optimize this
  identically to \passthrough{\lstinline!any\_of!} with a negated
  predicate.
\item
  \textbf{Example}:
  \passthrough{\lstinline!cpp     bool no\_zeros = std::none\_of(v.begin(), v.end(), [](int i)\{ return i == 0; \});!}
\end{itemize}

\hypertarget{stdfor_each}{%
\subsubsection{\texorpdfstring{4.
\texttt{std::for\_each}}{4. std::for\_each}}\label{stdfor_each}}

\begin{itemize}
\tightlist
\item
  \textbf{Analogy}: The ``Delivery Driver''. Stopping at every house to
  drop off a package.
\item
  \textbf{When to use it}: When you want to perform an action on every
  element (usually for side effects like logging or updating a hardware
  register).
\item
  \textbf{Complexity}: \(O(n)\).
\item
  \textbf{Godhood Tip}: Since C++17, use the execution policy
  \passthrough{\lstinline!std::execution::par!} to make this
  multi-threaded instantly!
\item
  \textbf{Hardware Sympathy}: Perfect for prefetching. The CPU sees the
  linear access pattern and starts pulling data into L1 cache before you
  even ask for it.
\item
  \textbf{Example}:
  \passthrough{\lstinline!cpp     std::for\_each(std::execution::par, v.begin(), v.end(), [](int\& i)\{ i *= 2; \});!}
\end{itemize}

\hypertarget{stdfind}{%
\subsubsection{\texorpdfstring{5.
\texttt{std::find}}{5. std::find}}\label{stdfind}}

\begin{itemize}
\tightlist
\item
  \textbf{Analogy}: ``Where's Waldo?''. Looking through a crowd until
  you find the exact match.
\item
  \textbf{When to use it}: Simple value searching in an unsorted
  container.
\item
  \textbf{Complexity}: \(O(n)\).
\item
  \textbf{Hardware Sympathy}: Linear scan. Extremely cache-friendly
  compared to \passthrough{\lstinline!std::set::find!} or
  \passthrough{\lstinline!std::map::find!}, which involve pointer
  chasing.
\item
  \textbf{Example}:
  \passthrough{\lstinline!cpp     auto it = std::find(v.begin(), v.end(), 42);!}
\end{itemize}

\hypertarget{stdfind_if}{%
\subsubsection{\texorpdfstring{6.
\texttt{std::find\_if}}{6. std::find\_if}}\label{stdfind_if}}

\begin{itemize}
\tightlist
\item
  \textbf{Analogy}: The ``Headhunter''. Looking for anyone who speaks 5
  languages and knows COBOL.
\item
  \textbf{When to use it}: Searching for an element that matches a
  specific, complex predicate.
\item
  \textbf{Complexity}: \(O(n)\).
\item
  \textbf{Common Pitfall}: Passing a heavy predicate by value. If the
  predicate has a large state, use a reference or a lambda.
\item
  \textbf{Example}:
  \passthrough{\lstinline!cpp     auto it = std::find\_if(v.begin(), v.end(), [](const auto\& emp)\{ return emp.salary > 200000; \});!}
\end{itemize}

\hypertarget{stdfind_if_not}{%
\subsubsection{\texorpdfstring{7.
\texttt{std::find\_if\_not}}{7. std::find\_if\_not}}\label{stdfind_if_not}}

\begin{itemize}
\tightlist
\item
  \textbf{Analogy}: The ``Odd One Out''. Looking for the first person
  who ISN'T wearing a uniform.
\item
  \textbf{When to use it}: Finding the first element that fails a
  condition.
\item
  \textbf{Complexity}: \(O(n)\).
\item
  \textbf{Example}:
  \passthrough{\lstinline!cpp     auto it = std::find\_if\_not(v.begin(), v.end(), [](int i)\{ return i \% 2 == 0; \});!}
\end{itemize}

\hypertarget{stdfind_end}{%
\subsubsection{\texorpdfstring{8.
\texttt{std::find\_end}}{8. std::find\_end}}\label{stdfind_end}}

\begin{itemize}
\tightlist
\item
  \textbf{Analogy}: The ``Last Occurrence''. Finding the \emph{last}
  time a specific sequence appeared in a stream.
\item
  \textbf{When to use it}: When you need the tail end of a sub-sequence
  (e.g., finding the last occurrence of a file extension).
\item
  \textbf{Complexity}: \(O(n \cdot m)\).
\item
  \textbf{Hardware Sympathy}: This is a heavy hitter. If the
  sub-sequence \(m\) is large, this can be slow. Consider C++17
  searchers for better performance.
\item
  \textbf{Example}:
  \passthrough{\lstinline!cpp     auto it = std::find\_end(text.begin(), text.end(), sub.begin(), sub.end());!}
\end{itemize}

\hypertarget{stdfind_first_of}{%
\subsubsection{\texorpdfstring{9.
\texttt{std::find\_first\_of}}{9. std::find\_first\_of}}\label{stdfind_first_of}}

\begin{itemize}
\tightlist
\item
  \textbf{Analogy}: The ``Scavenger Hunt''. Looking for any one of
  several target items.
\item
  \textbf{When to use it}: Searching for the first occurrence of any
  element from a set of values (e.g., finding the first punctuation mark
  in a string).
\item
  \textbf{Complexity}: \(O(n \cdot m)\).
\item
  \textbf{Example}:
  \passthrough{\lstinline"cpp     std::vector<char> delimiters = \{',', '.', ';', '!'\};\\n    auto it = std::find\_first\_of(str.begin(), str.end(), delimiters.begin(), delimiters.end());"}
\end{itemize}

\hypertarget{stdadjacent_find}{%
\subsubsection{\texorpdfstring{10.
\texttt{std::adjacent\_find}}{10. std::adjacent\_find}}\label{stdadjacent_find}}

\begin{itemize}
\tightlist
\item
  \textbf{Analogy}: The ``Glitch Spotter''. Finding two identical frames
  in a row in a video stream.
\item
  \textbf{When to use it}: Detecting duplicates that are positioned next
  to each other.
\item
  \textbf{Complexity}: \(O(n)\).
\item
  \textbf{Hardware Sympathy}: Only compares neighbors. High spatial
  locality.
\item
  \textbf{Example}:
  \passthrough{\lstinline!cpp     auto it = std::adjacent\_find(v.begin(), v.end());!}
\end{itemize}

\hypertarget{stdcount}{%
\subsubsection{\texorpdfstring{11.
\texttt{std::count}}{11. std::count}}\label{stdcount}}

\begin{itemize}
\tightlist
\item
  \textbf{Analogy}: The ``Census Taker''. Counting how many people named
  ``Smith'' live in the city.
\item
  \textbf{When to use it}: Simple frequency counting.
\item
  \textbf{Complexity}: \(O(n)\).
\item
  \textbf{Example}:
  \passthrough{\lstinline!cpp     int num\_sevens = std::count(v.begin(), v.end(), 7);!}
\end{itemize}

\hypertarget{stdcount_if}{%
\subsubsection{\texorpdfstring{12.
\texttt{std::count\_if}}{12. std::count\_if}}\label{stdcount_if}}

\begin{itemize}
\tightlist
\item
  \textbf{Analogy}: The ``Pollster''. Counting how many people plan to
  vote ``Yes''.
\item
  \textbf{When to use it}: Counting elements that match a dynamic
  condition.
\item
  \textbf{Complexity}: \(O(n)\).
\item
  \textbf{Godhood Tip}: Many modern compilers will auto-vectorize this
  using SIMD (Single Instruction Multiple Data) if the predicate is
  simple.
\item
  \textbf{Example}:
  \passthrough{\lstinline!cpp     int positives = std::count\_if(v.begin(), v.end(), [](int i)\{ return i > 0; \});!}
\end{itemize}

\hypertarget{stdmismatch}{%
\subsubsection{\texorpdfstring{13.
\texttt{std::mismatch}}{13. std::mismatch}}\label{stdmismatch}}

\begin{itemize}
\tightlist
\item
  \textbf{Analogy}: ``Spot the Difference''. Comparing two photos and
  finding the first pixel that changed.
\item
  \textbf{When to use it}: Comparing two sequences (e.g., two versions
  of a config file) to find where they diverge.
\item
  \textbf{Complexity}: \(O(n)\).
\item
  \textbf{Common Pitfall}: Ensure the second range is at least as long
  as the first, or use the C++14 four-iterator version to avoid
  out-of-bounds access.
\item
  \textbf{Example}:
  \passthrough{\lstinline!cpp     auto [it1, it2] = std::mismatch(v1.begin(), v1.end(), v2.begin(), v2.end());!}
\end{itemize}

\hypertarget{stdequal}{%
\subsubsection{\texorpdfstring{14.
\texttt{std::equal}}{14. std::equal}}\label{stdequal}}

\begin{itemize}
\tightlist
\item
  \textbf{Analogy}: The ``Clone Check''. Verifying two documents are
  identical.
\item
  \textbf{When to use it}: Deep comparison of two ranges.
\item
  \textbf{Complexity}: \(O(n)\).
\item
  \textbf{Hardware Sympathy}: Compilers often optimize this to
  \passthrough{\lstinline!memcmp!} for primitive types, which is the
  gold standard for performance.
\item
  \textbf{Example}:
  \passthrough{\lstinline!cpp     bool is\_equal = std::equal(v1.begin(), v1.end(), v2.begin());!}
\end{itemize}

\hypertarget{stdis_permutation}{%
\subsubsection{\texorpdfstring{15.
\texttt{std::is\_permutation}}{15. std::is\_permutation}}\label{stdis_permutation}}

\begin{itemize}
\tightlist
\item
  \textbf{Analogy}: The ``Anagram''. Checking if ``Listen'' and
  ``Silent'' have the same letters.
\item
  \textbf{When to use it}: Checking if two ranges have the same elements
  in any order.
\item
  \textbf{Complexity}: \(O(n^2)\) (Worst case).
\item
  \textbf{Godhood Tip}: This is expensive! If you need to do this often
  on large sets, sort both ranges first and use
  \passthrough{\lstinline!std::equal!} (\(O(n \log n)\) total).
\item
  \textbf{Example}:
  \passthrough{\lstinline!cpp     bool anagram = std::is\_permutation(word1.begin(), word1.end(), word2.begin());!}
\end{itemize}

\hypertarget{stdsearch}{%
\subsubsection{\texorpdfstring{16.
\texttt{std::search}}{16. std::search}}\label{stdsearch}}

\begin{itemize}
\tightlist
\item
  \textbf{Analogy}: ``Ctrl+F''. Searching for a specific word in a
  sentence.
\item
  \textbf{When to use it}: Finding a sub-sequence within a range.
\item
  \textbf{Complexity}: \(O(n \cdot m)\).
\item
  \textbf{Godhood Tip}: In C++17, you can pass a
  \passthrough{\lstinline!Searcher!} object (like
  \passthrough{\lstinline!std::boyer\_moore\_searcher!}) to achieve
  sub-linear performance (\(O(n/m)\)).
\item
  \textbf{Example}:
  \passthrough{\lstinline!cpp     auto it = std::search(text.begin(), text.end(), \\n                         std::boyer\_moore\_searcher(pattern.begin(), pattern.end()));!}
\end{itemize}

\hypertarget{stdsearch_n}{%
\subsubsection{\texorpdfstring{17.
\texttt{std::search\_n}}{17. std::search\_n}}\label{stdsearch_n}}

\begin{itemize}
\tightlist
\item
  \textbf{Analogy}: The ``Winning Streak''. Finding the first place
  where someone won 5 times in a row.
\item
  \textbf{When to use it}: Looking for \passthrough{\lstinline!n!}
  consecutive occurrences of a specific value.
\item
  \textbf{Complexity}: \(O(n)\).
\item
  \textbf{Example}:
  \passthrough{\lstinline!cpp     auto it = std::search\_n(v.begin(), v.end(), 5, 100); // 5 consecutive 100s!}
\end{itemize}

\hypertarget{stdcopy}{%
\subsubsection{\texorpdfstring{18.
\texttt{std::copy}}{18. std::copy}}\label{stdcopy}}

\begin{itemize}
\tightlist
\item
  \textbf{Analogy}: The ``Photocopier''. Making an exact duplicate of a
  stack of papers.
\item
  \textbf{When to use it}: Moving data from one range to another.
\item
  \textbf{Complexity}: \(O(n)\).
\item
  \textbf{Hardware Sympathy}: Usually compiles down to
  \passthrough{\lstinline!memcpy!}, the fastest possible way to move
  bytes in a computer.
\item
  \textbf{Example}:
  \passthrough{\lstinline!cpp     std::copy(src.begin(), src.end(), dest.begin());!}
\end{itemize}

\hypertarget{stdcopy_n}{%
\subsubsection{\texorpdfstring{19.
\texttt{std::copy\_n}}{19. std::copy\_n}}\label{stdcopy_n}}

\begin{itemize}
\tightlist
\item
  \textbf{Analogy}: The ``Limited Edition''. Only copying the first 10
  pages of a book.
\item
  \textbf{When to use it}: When you know exactly how many elements to
  move, avoiding the need for an end iterator.
\item
  \textbf{Complexity}: \(O(n)\).
\item
  \textbf{Example}:
  \passthrough{\lstinline!cpp     std::copy\_n(src.begin(), 10, dest.begin());!}
\end{itemize}

\hypertarget{stdcopy_if}{%
\subsubsection{\texorpdfstring{20.
\texttt{std::copy\_if}}{20. std::copy\_if}}\label{stdcopy_if}}

\begin{itemize}
\tightlist
\item
  \textbf{Analogy}: The ``Filter''. Only copying the ``VIP'' names from
  the guest list.
\item
  \textbf{When to use it}: Moving data that meets a certain criteria.
\item
  \textbf{Complexity}: \(O(n)\).
\item
  \textbf{Hardware Sympathy}: Can be slower than
  \passthrough{\lstinline!std::copy!} due to branch mispredictions in
  the predicate if the data is random.
\item
  \textbf{Example}:
  \passthrough{\lstinline!cpp     std::copy\_if(src.begin(), src.end(), std::back\_inserter(dest), [](int i)\{ return i > 0; \});!}
\end{itemize}

\hypertarget{stdcopy_backward}{%
\subsubsection{\texorpdfstring{21.
\texttt{std::copy\_backward}}{21. std::copy\_backward}}\label{stdcopy_backward}}

\begin{itemize}
\tightlist
\item
  \textbf{Analogy}: The ``Reverse Conveyor''. Copying items but starting
  from the end of the destination to avoid overwriting.
\item
  \textbf{When to use it}: When source and destination ranges overlap
  and the destination is further ahead in memory (shifting right).
\item
  \textbf{Complexity}: \(O(n)\).
\item
  \textbf{Example}:
  \passthrough{\lstinline!cpp     std::copy\_backward(v.begin(), v.begin() + 5, v.begin() + 7);!}
\end{itemize}

\hypertarget{stdmove-algorithm}{%
\subsubsection{\texorpdfstring{22. \texttt{std::move}
(algorithm)}{22. std::move (algorithm)}}\label{stdmove-algorithm}}

\begin{itemize}
\tightlist
\item
  \textbf{Analogy}: The ``Moving Van''. Not just copying, but actually
  taking the furniture out of the old house.
\item
  \textbf{When to use it}: Efficiency! Use when you don't need the
  source elements anymore and they support move semantics.
\item
  \textbf{Complexity}: \(O(n)\).
\item
  \textbf{Hardware Sympathy}: For types like
  \passthrough{\lstinline!std::string!} or
  \passthrough{\lstinline!std::vector!}, this is vastly faster than
  \passthrough{\lstinline!copy!} because it just swaps pointers.
\item
  \textbf{Example}:
  \passthrough{\lstinline!cpp     std::move(src.begin(), src.end(), dest.begin());!}
\end{itemize}

\hypertarget{stdmove_backward}{%
\subsubsection{\texorpdfstring{23.
\texttt{std::move\_backward}}{23. std::move\_backward}}\label{stdmove_backward}}

\begin{itemize}
\tightlist
\item
  \textbf{Analogy}: Shifting a row of expensive vases to the right,
  moving the last one first to avoid breakage.
\item
  \textbf{When to use it}: Overlapping ranges where the destination
  starts inside the source range and is to the right.
\item
  \textbf{Complexity}: \(O(n)\).
\item
  \textbf{Example}:
  \passthrough{\lstinline!cpp     std::move\_backward(src.begin(), src.end(), dest.end());!}
\end{itemize}

\hypertarget{stdswap_ranges}{%
\subsubsection{\texorpdfstring{24.
\texttt{std::swap\_ranges}}{24. std::swap\_ranges}}\label{stdswap_ranges}}

\begin{itemize}
\tightlist
\item
  \textbf{Analogy}: The ``Trading Places''. Two rows of students
  swapping seats with each other simultaneously.
\item
  \textbf{When to use it}: Swapping chunks of data between containers
  without temporary allocations.
\item
  \textbf{Complexity}: \(O(n)\).
\item
  \textbf{Example}:
  \passthrough{\lstinline!cpp     std::swap\_ranges(v1.begin(), v1.end(), v2.begin());!}
\end{itemize}

\hypertarget{stdtransform}{%
\subsubsection{\texorpdfstring{25.
\texttt{std::transform}}{25. std::transform}}\label{stdtransform}}

\begin{itemize}
\tightlist
\item
  \textbf{Analogy}: The ``Assembly Line''. Every part comes in raw and
  gets polished on its way out.
\item
  \textbf{When to use it}: Applying a function to every element and
  storing the result elsewhere. This is the ``Map'' in MapReduce.
\item
  \textbf{Complexity}: \(O(n)\).
\item
  \textbf{Godhood Tip}: You can also use the binary version to combine
  two ranges into one (e.g., adding two vectors).
\item
  \textbf{Example}:
  \passthrough{\lstinline!cpp     std::transform(v.begin(), v.end(), v.begin(), [](int i)\{ return i * i; \});!}
\end{itemize}

\hypertarget{stdreplace}{%
\subsubsection{\texorpdfstring{26.
\texttt{std::replace}}{26. std::replace}}\label{stdreplace}}

\begin{itemize}
\tightlist
\item
  \textbf{Analogy}: ``Search and Replace''. Changing every ``Apple'' to
  ``Orange'' in a document.
\item
  \textbf{When to use it}: Simple value replacement across a container.
\item
  \textbf{Complexity}: \(O(n)\).
\item
  \textbf{Example}:
  \passthrough{\lstinline!cpp     std::replace(v.begin(), v.end(), old\_val, new\_val);!}
\end{itemize}

\hypertarget{stdreplace_if}{%
\subsubsection{\texorpdfstring{27.
\texttt{std::replace\_if}}{27. std::replace\_if}}\label{stdreplace_if}}

\begin{itemize}
\tightlist
\item
  \textbf{Analogy}: The ``Tax Man''. Replacing every salary over 100k
  with a fixed ``Cap''.
\item
  \textbf{When to use it}: Conditional replacement based on a predicate.
\item
  \textbf{Complexity}: \(O(n)\).
\item
  \textbf{Example}:
  \passthrough{\lstinline!cpp     std::replace\_if(v.begin(), v.end(), [](int i)\{ return i > 100; \}, 100);!}
\end{itemize}

\hypertarget{stdfill}{%
\subsubsection{\texorpdfstring{28.
\texttt{std::fill}}{28. std::fill}}\label{stdfill}}

\begin{itemize}
\tightlist
\item
  \textbf{Analogy}: The ``Paint Bucket''. Filling a whole canvas with a
  single color.
\item
  \textbf{When to use it}: Initializing a range (e.g., a buffer) with a
  constant value.
\item
  \textbf{Complexity}: \(O(n)\).
\item
  \textbf{Hardware Sympathy}: Compiles to
  \passthrough{\lstinline!memset!} for byte-sized types. The fastest
  possible memory write.
\item
  \textbf{Example}:
  \passthrough{\lstinline!cpp     std::fill(v.begin(), v.end(), 0);!}
\end{itemize}

\hypertarget{stdfill_n}{%
\subsubsection{\texorpdfstring{29.
\texttt{std::fill\_n}}{29. std::fill\_n}}\label{stdfill_n}}

\begin{itemize}
\tightlist
\item
  \textbf{Analogy}: ``First 10 are Free''. Only painting the first 10
  items in a row.
\item
  \textbf{When to use it}: When you have a pointer/iterator and a count
  but no end iterator.
\item
  \textbf{Complexity}: \(O(n)\).
\item
  \textbf{Example}:
  \passthrough{\lstinline!cpp     std::fill\_n(v.begin(), 10, -1);!}
\end{itemize}

\hypertarget{stdgenerate}{%
\subsubsection{\texorpdfstring{30.
\texttt{std::generate}}{30. std::generate}}\label{stdgenerate}}

\begin{itemize}
\tightlist
\item
  \textbf{Analogy}: The ``Random Number Generator''. Calling a function
  to create a new value for every slot.
\item
  \textbf{When to use it}: Filling a range with dynamic or random
  values.
\item
  \textbf{Complexity}: \(O(n)\).
\item
  \textbf{Example}:
  \passthrough{\lstinline!cpp     std::generate(v.begin(), v.end(), std::rand);!}
\end{itemize}

\hypertarget{stdgenerate_n}{%
\subsubsection{\texorpdfstring{31.
\texttt{std::generate\_n}}{31. std::generate\_n}}\label{stdgenerate_n}}

\begin{itemize}
\tightlist
\item
  \textbf{Analogy}: ``Print 5 Tickets''. Generating a specific number of
  new items.
\item
  \textbf{When to use it}: Populating a specific count of elements
  dynamically into a container.
\item
  \textbf{Complexity}: \(O(n)\).
\item
  \textbf{Example}:
  \passthrough{\lstinline!cpp     std::generate\_n(std::back\_inserter(v), 5, []()\{ return rand() \% 100; \});!}
\end{itemize}

\hypertarget{stdremove-the-shifting-seats}{%
\subsubsection{\texorpdfstring{32. \texttt{std::remove} (The Shifting
Seats)}{32. std::remove (The Shifting Seats)}}\label{stdremove-the-shifting-seats}}

\begin{itemize}
\tightlist
\item
  \textbf{Analogy}: Moving all the ``Empty'' chairs to the back of the
  room so the front rows are full and usable.
\item
  \textbf{When to use it}: ``Deleting'' elements from a container
  (Vector/Array).
\item
  \textbf{Complexity}: \(O(n)\).
\item
  \textbf{CRITICAL WARNING}: It doesn't actually change the size of the
  container! You MUST use the \textbf{Erase-Remove Idiom}.
\item
  \textbf{Hardware Sympathy}: Very fast because it only performs
  \(O(n)\) moves instead of \(O(n^2)\) shifts.
\item
  \textbf{Example}:
  \passthrough{\lstinline!cpp     v.erase(std::remove(v.begin(), v.end(), 99), v.end());!}
\end{itemize}

\hypertarget{stdremove_if}{%
\subsubsection{\texorpdfstring{33.
\texttt{std::remove\_if}}{33. std::remove\_if}}\label{stdremove_if}}

\begin{itemize}
\tightlist
\item
  \textbf{Analogy}: ``Excommunicated''. Shifting everyone who failed a
  test to the back of the line.
\item
  \textbf{When to use it}: Conditional removal from a collection.
\item
  \textbf{Complexity}: \(O(n)\).
\item
  \textbf{Example}:
  \passthrough{\lstinline!cpp     v.erase(std::remove\_if(v.begin(), v.end(), [](int i)\{ return i < 0; \}), v.end());!}
\end{itemize}

\hypertarget{stdremove_copy}{%
\subsubsection{\texorpdfstring{34.
\texttt{std::remove\_copy}}{34. std::remove\_copy}}\label{stdremove_copy}}

\begin{itemize}
\tightlist
\item
  \textbf{Analogy}: ``Selective Copying''. Copying a list but skipping
  specific ``Banned'' names.
\item
  \textbf{When to use it}: When you want to keep the original data but
  need a cleaned-up copy.
\item
  \textbf{Complexity}: \(O(n)\).
\item
  \textbf{Example}:
  \passthrough{\lstinline!cpp     std::remove\_copy(src.begin(), src.end(), std::back\_inserter(dest), 99);!}
\end{itemize}

\hypertarget{stdremove_copy_if}{%
\subsubsection{\texorpdfstring{35.
\texttt{std::remove\_copy\_if}}{35. std::remove\_copy\_if}}\label{stdremove_copy_if}}

\begin{itemize}
\tightlist
\item
  \textbf{Analogy}: ``The Purge''. Copying a list but leaving out anyone
  who doesn't meet the criteria.
\item
  \textbf{When to use it}: Copying only elements that fail a specific
  condition.
\item
  \textbf{Complexity}: \(O(n)\).
\item
  \textbf{Example}:
  \passthrough{\lstinline!cpp     std::remove\_copy\_if(src.begin(), src.end(), std::back\_inserter(dest), [](int i)\{ return i < 10; \});!}
\end{itemize}

\hypertarget{stdunique}{%
\subsubsection{\texorpdfstring{36.
\texttt{std::unique}}{36. std::unique}}\label{stdunique}}

\begin{itemize}
\tightlist
\item
  \textbf{Analogy}: ``Stop Repeating Yourself!''. If someone says the
  same word twice in a row, tell them to stop.
\item
  \textbf{When to use it}: Removing \emph{consecutive} duplicates.
\item
  \textbf{Complexity}: \(O(n)\).
\item
  \textbf{Godhood Tip}: To remove \emph{all} duplicates, you must
  \passthrough{\lstinline!sort()!} before calling
  \passthrough{\lstinline!unique()!}.
\item
  \textbf{Example}:
  \passthrough{\lstinline!cpp     v.erase(std::unique(v.begin(), v.end()), v.end());!}
\end{itemize}

\hypertarget{stdunique_copy}{%
\subsubsection{\texorpdfstring{37.
\texttt{std::unique\_copy}}{37. std::unique\_copy}}\label{stdunique_copy}}

\begin{itemize}
\tightlist
\item
  \textbf{Analogy}: ``Recording the Highlights''. Copying a sequence but
  only taking one instance of any consecutive group.
\item
  \textbf{When to use it}: Creating a ``de-duplicated'' version of a
  range.
\item
  \textbf{Complexity}: \(O(n)\).
\item
  \textbf{Example}:
  \passthrough{\lstinline!cpp     std::unique\_copy(src.begin(), src.end(), std::back\_inserter(dest));!}
\end{itemize}

\hypertarget{stdreverse}{%
\subsubsection{\texorpdfstring{38.
\texttt{std::reverse}}{38. std::reverse}}\label{stdreverse}}

\begin{itemize}
\tightlist
\item
  \textbf{Analogy}: ``The Rewind''. Flipping the whole sequence upside
  down.
\item
  \textbf{When to use it}: When you need the order completely inverted
  (e.g., converting big-endian to little-endian manually).
\item
  \textbf{Complexity}: \(O(n)\).
\item
  \textbf{Hardware Sympathy}: High bandwidth usage. Swaps elements from
  outside-in, moving linearly through memory.
\item
  \textbf{Example}:
  \passthrough{\lstinline!cpp     std::reverse(v.begin(), v.end());!}
\end{itemize}

\hypertarget{stdreverse_copy}{%
\subsubsection{\texorpdfstring{39.
\texttt{std::reverse\_copy}}{39. std::reverse\_copy}}\label{stdreverse_copy}}

\begin{itemize}
\tightlist
\item
  \textbf{Analogy}: ``Mirror Image''. Copying a list into another
  container but in reverse order.
\item
  \textbf{When to use it}: Keeping the original order while obtaining a
  reversed version.
\item
  \textbf{Complexity}: \(O(n)\).
\item
  \textbf{Example}:
  \passthrough{\lstinline!cpp     std::reverse\_copy(src.begin(), src.end(), dest.begin());!}
\end{itemize}

\hypertarget{stdrotate-the-pivot-dance}{%
\subsubsection{\texorpdfstring{40. \texttt{std::rotate} (The Pivot
Dance)}{40. std::rotate (The Pivot Dance)}}\label{stdrotate-the-pivot-dance}}

\begin{itemize}
\tightlist
\item
  \textbf{Analogy}: ``The Conveyor Belt''. Moving the 3rd item to the
  front and shifting everything else.
\item
  \textbf{When to use it}: Cyclic shifts. This is the magic behind
  moving an element from index A to index B.
\item
  \textbf{Complexity}: \(O(n)\).
\item
  \textbf{Hardware Sympathy}: One of the most highly optimized
  algorithms. Can be used for ``O(1)'' front-deletion in a vector if
  order doesn't matter (rotate + pop\_back).
\item
  \textbf{Example}:
  \passthrough{\lstinline!cpp     std::rotate(v.begin(), v.begin() + 2, v.end());!}
\end{itemize}

\hypertarget{stdrotate_copy}{%
\subsubsection{\texorpdfstring{41.
\texttt{std::rotate\_copy}}{41. std::rotate\_copy}}\label{stdrotate_copy}}

\begin{itemize}
\tightlist
\item
  \textbf{Analogy}: ``Circular Snapshot''. Taking a picture of the
  conveyor belt after it has rotated.
\item
  \textbf{When to use it}: Getting a shifted copy of a range without
  modifying the original.
\item
  \textbf{Complexity}: \(O(n)\).
\item
  \textbf{Example}:
  \passthrough{\lstinline!cpp     std::rotate\_copy(src.begin(), src.begin() + 3, src.end(), dest.begin());!}
\end{itemize}

\hypertarget{stdshift_left-c20}{%
\subsubsection{\texorpdfstring{42. \texttt{std::shift\_left}
(C++20)}{42. std::shift\_left (C++20)}}\label{stdshift_left-c20}}

\begin{itemize}
\tightlist
\item
  \textbf{Analogy}: ``The Slide''. Everyone slides to the left by 2
  seats. The people at the far left are discarded.
\item
  \textbf{When to use it}: Shifting data without the overhead of
  rotation.
\item
  \textbf{Complexity}: \(O(n)\).
\item
  \textbf{Hardware Sympathy}: Efficiently moves data without wrapping
  around, useful in buffer management.
\item
  \passthrough{\lstinline!cpp     std::shift\_left(v.begin(), v.end(), 2);!}
\end{itemize}

\hypertarget{stdshift_right-c20}{%
\subsubsection{\texorpdfstring{43. \texttt{std::shift\_right}
(C++20)}{43. std::shift\_right (C++20)}}\label{stdshift_right-c20}}

\begin{itemize}
\tightlist
\item
  \textbf{Analogy}: ``The Push''. Everyone moves right. The people at
  the end are pushed out of the room.
\item
  \textbf{When to use it}: Shifting data right.
\item
  \textbf{Complexity}: \(O(n)\).
\item
  \passthrough{\lstinline!cpp     std::shift\_right(v.begin(), v.end(), 2);!}
\end{itemize}

\hypertarget{stdshuffle}{%
\subsubsection{\texorpdfstring{44.
\texttt{std::shuffle}}{44. std::shuffle}}\label{stdshuffle}}

\begin{itemize}
\tightlist
\item
  \textbf{Analogy}: ``The Vegas Dealer''. Mixing the deck so perfectly
  that the outcome is statistically unpredictable.
\item
  \textbf{When to use it}: Randomizing a range for simulations or games.
\item
  \textbf{Complexity}: \(O(n)\).
\item
  \textbf{Common Pitfall}: Using \passthrough{\lstinline!rand()!} or
  \passthrough{\lstinline!random\_shuffle!} (which are deprecated/poor).
  Use \passthrough{\lstinline!std::mt19937!} for true randomness.
\item
  \textbf{Example}:
  \passthrough{\lstinline!cpp     std::shuffle(v.begin(), v.end(), std::mt19937\{std::random\_device\{\}()\});!}
\end{itemize}

\hypertarget{stdis_partitioned}{%
\subsubsection{\texorpdfstring{45.
\texttt{std::is\_partitioned}}{45. std::is\_partitioned}}\label{stdis_partitioned}}

\begin{itemize}
\tightlist
\item
  \textbf{Analogy}: ``Sorted by Side''. Checking if all the ``Blue''
  shirts are on the left and ``Red'' shirts on the right.
\item
  \textbf{When to use it}: Validating if a range has been successfully
  divided by a predicate.
\item
  \textbf{Complexity}: \(O(n)\).
\item
  \textbf{Example}:
  \passthrough{\lstinline!cpp     bool is\_p = std::is\_partitioned(v.begin(), v.end(), [](int i)\{ return i < 0; \});!}
\end{itemize}

\hypertarget{stdpartition}{%
\subsubsection{\texorpdfstring{46.
\texttt{std::partition}}{46. std::partition}}\label{stdpartition}}

\begin{itemize}
\tightlist
\item
  \textbf{Analogy}: ``The Middle School Dance''. Boys on the left, girls
  on the right.
\item
  \textbf{When to use it}: Fast separation of elements based on a
  condition without the cost of a full sort.
\item
  \textbf{Complexity}: \(O(n)\).
\item
  \textbf{Hardware Sympathy}: Much faster than
  \passthrough{\lstinline!std::sort!}. This is the fundamental building
  block of QuickSort.
\item
  \textbf{Example}:
  \passthrough{\lstinline!cpp     auto it = std::partition(v.begin(), v.end(), [](int i)\{ return i \% 2 == 0; \});!}
\end{itemize}

\hypertarget{stdstable_partition}{%
\subsubsection{\texorpdfstring{47.
\texttt{std::stable\_partition}}{47. std::stable\_partition}}\label{stdstable_partition}}

\begin{itemize}
\tightlist
\item
  \textbf{Analogy}: ``The Dance (Respecting Friendships)''. Boys on the
  left, girls on the right, but everyone keeps their original relative
  order with their friends.
\item
  \textbf{When to use it}: When the relative order of elements within
  the two partitions must be preserved.
\item
  \textbf{Complexity}: \(O(n \log n)\) or \(O(n)\) if extra memory is
  available.
\item
  \textbf{Example}:
  \passthrough{\lstinline!cpp     std::stable\_partition(v.begin(), v.end(), [](int i)\{ return i > 100; \});!}
\end{itemize}

\hypertarget{stdpartition_copy}{%
\subsubsection{\texorpdfstring{48.
\texttt{std::partition\_copy}}{48. std::partition\_copy}}\label{stdpartition_copy}}

\begin{itemize}
\tightlist
\item
  \textbf{Analogy}: ``Sorting the Mail''. Putting ``Bills'' in one bin
  and ``Letters'' in another.
\item
  \textbf{When to use it}: Moving elements into two separate containers
  based on a boolean condition.
\item
  \textbf{Complexity}: \(O(n)\).
\item
  \textbf{Example}:
  \passthrough{\lstinline!cpp     std::partition\_copy(src.begin(), src.end(), std::back\_inserter(v1), std::back\_inserter(v2), pred);!}
\end{itemize}

\hypertarget{stdpartition_point}{%
\subsubsection{\texorpdfstring{49.
\texttt{std::partition\_point}}{49. std::partition\_point}}\label{stdpartition_point}}

\begin{itemize}
\tightlist
\item
  \textbf{Analogy}: ``The Boundary Line''. Finding the exact spot where
  the ``Blue'' shirts end and ``Red'' shirts begin.
\item
  \textbf{When to use it}: Finding the pivot iterator in a previously
  partitioned range.
\item
  \textbf{Complexity}: \(O(\log n)\).
\item
  \textbf{Example}:
  \passthrough{\lstinline!cpp     auto it = std::partition\_point(v.begin(), v.end(), [](int i)\{ return i < 10; \});!}
\end{itemize}

\hypertarget{stdsort}{%
\subsubsection{\texorpdfstring{50.
\texttt{std::sort}}{50. std::sort}}\label{stdsort}}

\begin{itemize}
\tightlist
\item
  \textbf{Analogy}: ``The Library''. Putting every book in perfect
  alphabetical order.
\item
  \textbf{When to use it}: General purpose sorting.
\item
  \textbf{Complexity}: \(O(n \log n)\).
\item
  \textbf{Hardware Sympathy}: Typically implements \textbf{Introsort}
  (QuickSort + HeapSort + InsertionSort). It is extremely cache-friendly
  and avoids the \(O(n^2)\) worst-case of pure QuickSort.
\item
  \textbf{Example}:
  \passthrough{\lstinline!cpp     std::sort(v.begin(), v.end());!}
\end{itemize}

\hypertarget{stdstable_sort}{%
\subsubsection{\texorpdfstring{51.
\texttt{std::stable\_sort}}{51. std::stable\_sort}}\label{stdstable_sort}}

\begin{itemize}
\tightlist
\item
  \textbf{Analogy}: ``Sorting by Group, then Rank''. Sorting students by
  score, but ensuring students with the same score stay in the order
  they were in.
\item
  \textbf{When to use it}: When relative order of ``equal'' elements is
  semantically important.
\item
  \textbf{Complexity}: \(O(n \log^2 n)\) (or \(O(n \log n)\) with extra
  memory).
\item
  \textbf{Example}:
  \passthrough{\lstinline!cpp     std::stable\_sort(v.begin(), v.end());!}
\end{itemize}

\hypertarget{stdpartial_sort}{%
\subsubsection{\texorpdfstring{52.
\texttt{std::partial\_sort}}{52. std::partial\_sort}}\label{stdpartial_sort}}

\begin{itemize}
\tightlist
\item
  \textbf{Analogy}: ``The Top 10 Leaderboard''. Finding the 10 fastest
  runners and sorting them, while the other 990 stay in any order.
\item
  \textbf{When to use it}: When you only need the smallest/largest \(k\)
  elements in order.
\item
  \textbf{Complexity}: \(O(n \log k)\).
\item
  \textbf{Hardware Sympathy}: Uses a heap internally. Much faster than a
  full sort if \(k \ll n\).
\item
  \textbf{Example}:
  \passthrough{\lstinline!cpp     std::partial\_sort(v.begin(), v.begin() + 5, v.end()); // Top 5 are sorted!}
\end{itemize}

\hypertarget{stdpartial_sort_copy}{%
\subsubsection{\texorpdfstring{53.
\texttt{std::partial\_sort\_copy}}{53. std::partial\_sort\_copy}}\label{stdpartial_sort_copy}}

\begin{itemize}
\tightlist
\item
  \textbf{Analogy}: ``Extracting the Top 10''. Finding the 10 best items
  and copying them to a separate list, sorted.
\item
  \textbf{When to use it}: Creating leaderboards without modifying the
  original dataset.
\item
  \textbf{Complexity}: \(O(n \log k)\).
\item
  \passthrough{\lstinline!cpp     std::partial\_sort\_copy(src.begin(), src.end(), top\_5.begin(), top\_5.end());!}
\end{itemize}

\hypertarget{stdis_sorted}{%
\subsubsection{\texorpdfstring{54.
\texttt{std::is\_sorted}}{54. std::is\_sorted}}\label{stdis_sorted}}

\begin{itemize}
\tightlist
\item
  \textbf{Analogy}: ``The Quality Control Check''. Making sure every
  single item on the conveyor belt is in the correct order.
\item
  \textbf{When to use it}: Verification and assertions in
  high-reliability code.
\item
  \textbf{Complexity}: \(O(n)\).
\item
  \passthrough{\lstinline!cpp     assert(std::is\_sorted(v.begin(), v.end()));!}
\end{itemize}

\hypertarget{stdis_sorted_until}{%
\subsubsection{\texorpdfstring{55.
\texttt{std::is\_sorted\_until}}{55. std::is\_sorted\_until}}\label{stdis_sorted_until}}

\begin{itemize}
\tightlist
\item
  \textbf{Analogy}: ``Finding the Point of Failure''. Finding the first
  item that breaks the sorted order.
\item
  \textbf{When to use it}: Identifying how much of a prefix is already
  sorted.
\item
  \textbf{Complexity}: \(O(n)\).
\item
  \passthrough{\lstinline!cpp     auto it = std::is\_sorted\_until(v.begin(), v.end());!}
\end{itemize}

\hypertarget{stdnth_element-the-median-finder}{%
\subsubsection{\texorpdfstring{56. \texttt{std::nth\_element} (The
Median
Finder)}{56. std::nth\_element (The Median Finder)}}\label{stdnth_element-the-median-finder}}

\begin{itemize}
\tightlist
\item
  \textbf{Analogy}: ``Finding the Middle Person''. Putting the person of
  median height in the center, and ensuring everyone shorter is to their
  left.
\item
  \textbf{When to use it}: Finding the median, the 99th percentile, or
  the top \(k\)-th element without the cost of sorting.
\item
  \textbf{Complexity}: \(O(n)\) (Average).
\item
  \textbf{Godhood Tip}: This is arguably the most under-used powerful
  algorithm in the STL. It's essentially a partial QuickSort.
\item
  \passthrough{\lstinline!cpp     std::nth\_element(v.begin(), v.begin() + v.size()/2, v.end());!}
\end{itemize}

\hypertarget{stdlower_bound}{%
\subsubsection{\texorpdfstring{57.
\texttt{std::lower\_bound}}{57. std::lower\_bound}}\label{stdlower_bound}}

\begin{itemize}
\tightlist
\item
  \textbf{Analogy}: ``The Insertion Point''. Finding the first place you
  could insert a value without breaking the sorted order.
\item
  \textbf{When to use it}: Binary search for the first element \(\ge\)
  value.
\item
  \textbf{Complexity}: \(O(\log n)\).
\item
  \textbf{CRITICAL}: The range MUST be sorted.
\item
  \textbf{Hardware Sympathy}: While \(O(\log n)\) is fast, large jumps
  in binary search can cause cache misses. For small ranges, a linear
  search (\passthrough{\lstinline!std::find!}) can actually be faster.
\item
  \textbf{Example}:
  \passthrough{\lstinline!cpp     auto it = std::lower\_bound(v.begin(), v.end(), 42);!}
\end{itemize}

\hypertarget{stdupper_bound}{%
\subsubsection{\texorpdfstring{58.
\texttt{std::upper\_bound}}{58. std::upper\_bound}}\label{stdupper_bound}}

\begin{itemize}
\tightlist
\item
  \textbf{Analogy}: ``The Last Insertion Point''. Finding the last
  possible place to insert a value.
\item
  \textbf{When to use it}: Binary search for the first element \(>\)
  value.
\item
  \textbf{Complexity}: \(O(\log n)\).
\item
  \textbf{Example}:
  \passthrough{\lstinline!cpp     auto it = std::upper\_bound(v.begin(), v.end(), 42);!}
\end{itemize}

\hypertarget{stdequal_range}{%
\subsubsection{\texorpdfstring{59.
\texttt{std::equal\_range}}{59. std::equal\_range}}\label{stdequal_range}}

\begin{itemize}
\tightlist
\item
  \textbf{Analogy}: ``The Target Zone''. Finding the beginning and end
  of all instances of a specific value.
\item
  \textbf{When to use it}: When you need both the first and last
  position of a value in a sorted range.
\item
  \textbf{Complexity}: \(O(\log n)\).
\item
  \passthrough{\lstinline!cpp     auto [first, last] = std::equal\_range(v.begin(), v.end(), 42);!}
\end{itemize}

\hypertarget{stdbinary_search}{%
\subsubsection{\texorpdfstring{60.
\texttt{std::binary\_search}}{60. std::binary\_search}}\label{stdbinary_search}}

\begin{itemize}
\tightlist
\item
  \textbf{Analogy}: ``The Yes/No Question''. Checking if a book is in
  the library without checking how many copies there are.
\item
  \textbf{When to use it}: Existence check in a sorted range when you
  don't need the iterator.
\item
  \textbf{Complexity}: \(O(\log n)\).
\item
  \passthrough{\lstinline!cpp     bool exists = std::binary\_search(v.begin(), v.end(), 42);!}
\end{itemize}

---\n
